% ============================================
% 第10章：人的判断权：AI越强，人越不能退后
% 悬索桥 · 判断权
% ============================================

\chapter{人的判断权：AI越强，人越不能退后}

\vspace{0.5cm}

\begin{center}
{\small\color{muted} 悬索桥跨越最长的距离。}
{\small\color{muted} 但它依赖最深的锚——锚碇，埋在地下，看不见，但绝对不能松。}
{\small\color{muted} 人的判断权，就是教育的锚碇。}
\end{center}

\vspace{1cm}

\section{``Skill不是取经路''}

2026年6月3日，师宇在知识星球上写了一篇博文，标题叫``Skill不是取经路''。里面有一段话，我反复读了很多遍：

\begin{shiyu-inline}
``Skill不是取经路。把经历写成'遇到妖怪就打、打不过找佛祖'，并不会让另一只猴子成为孙悟空。真正的能力不是知道规则，而是知道何时打、何时收手、何时求助、求谁，以及为什么还要往前走。''

\hfill ——师宇博文，2026年6月3日
\end{shiyu-inline}

这段话击中了AI时代最核心的问题。

Skill——AI技能包——是我这两年一直在推广的东西。把专家的隐性经验打包成可调用的流程：做什么、怎么做、要注意什么、什么情况下不能这么做。它让一个新手能借助AI，做出接近专家水平的东西。

但师宇在这篇博文里问了一个更尖锐的问题：\textbf{有了Skill，这个人就真的成了专家吗？}

孙悟空的取经经验，如果写成``遇到妖怪就打，打不过找佛祖''，当然也对。但另一只猴子读了这段话，不会变成孙悟空。因为它不知道\textbf{什么时候}该打，什么时候不该打，什么时候该求人，求谁，以及——\textbf{为什么非去不可。}

\section{AI把``做出来''变容易，把``判断''变难}

师宇在同一篇博文里继续写道：

\begin{shiyu-inline}
``AI把'做出来'的门槛降得很低，却把'知道好不好'的门槛抬得更高。Skill可能拉平动作表层，却拉远判断深层。它给普通人专家的外壳，也让真正专家更快迭代。''

\hfill ——师宇博文，2026年6月3日
\end{shiyu-inline}

\textbf{Skill给普通人专家的外壳。}

这句话让我想了很久。一个学生，用我做好的桥梁Skill，输入几个参数，AI生成了一座漂亮的斜拉桥模型。看起来，他``会''了。但如果你问他：``为什么塔高取这个值？为什么索距取这个值？如果地基条件变了，这个设计还能用吗？''——他答不上来。

\textbf{他有的不是能力。他有的是一套能生成专家外壳的工具。}

\begin{insight}
\label{insight:chap10}

\textbf{真正危险的，不是交出执行，是交出判断。}

AI可以帮你写。可以帮你画。可以帮你算。可以帮你归纳你是什么人。可以生成一套``如何对待你''的策略。它越强，人越容易舒服地退到后面。

但真正该被训练的，恰恰是\textbf{人不退后的能力}。人要知道自己为什么做这件事。知道什么值得做。知道哪里不够好。知道何时该慢下来。知道什么不能交出去。
\end{insight}

\section{``人与人的东西''仍被认为最重要}

2026年7月，我在教学理念综合分析中写下了这样一段话：

\begin{shiyu-inline}
``即使平台高度智能化，'人与人的东西'仍被认为最重要。材料一方面追求动态问题、即时排名和AI汇总，另一方面又明确说积分只是辅助、线下活动更重要，并建议学生的迭代过程应由学生自己总结，而不是交给AI。真正成熟的AI教学，可能不在于自动化更多，而在于知道哪些认知劳动必须保留给人。''

\hfill ——教学理念综合分析，2026年7月
\end{shiyu-inline}

这段话是我对两年教学实验的\textbf{自我限制}。

我的平台可以做很多事情。它可以自动评分、自动汇总反馈、自动生成知识图谱、自动检测学生的盲区。但它不能替代一件事：\textbf{学生自己总结自己的迭代过程。}

为什么？因为总结，是\textbf{判断}的核心动作。你回头看自己做了什么，哪些对了，哪些错了，为什么错，下次怎么改——这个动作，AI不能替你做。不是因为它做不到。是因为\textbf{它做了，就不是你的了。}

\section{交出判断的那一刻}

2026年5月31日，师宇写了一篇短文，标题叫``扫描我的电脑，看看你能做点什么？''。他写道：

\begin{shiyu-inline}
``这个提示词不是人告诉AI做什么，而是把判断权交给AI，让AI想想能为人做什么。这当然很强，也确实很诱人。可一旦人习惯了把判断权交出去，所谓赋能就会滑向献祭。''

\hfill ——师宇博文，2026年5月31日
\end{shiyu-inline}

\textbf{赋能滑向献祭。}

这是全书最让我不安的四个字。AI帮你做越来越多的事，你越来越轻松，越来越依赖它。有一天，你发现你不再需要思考了。AI帮你思考。你只需要点头。你不需要判断什么值得做——AI告诉你。你不需要知道哪里不够好——AI帮你改好了。你不需要决定什么不能交出去——因为\textbf{你已经不记得什么东西是重要的了}。

这就是悬索桥的锚碇。

\section{悬索桥的隐喻}

悬索桥跨越最长的距离。但它依赖最深的锚——锚碇，埋在地下，看不见，但绝对不能松。锚碇松了，整座桥就塌了。

\textbf{人的判断权，就是教育的锚碇。}

AI可以帮你搭桥。可以帮你画图。可以帮你写报告。可以帮你分析数据。但它不能帮你判断：这座桥对不对？这张图哪里不对？这个报告有没有注入真实经验？这个数据背后是什么物理现象？

这些判断，必须由人来做。不是因为AI做不了。是因为\textbf{AI做了，人就不再会做了}。而人不再会判断的时候，那座桥就会塌——不是物理上的塌，是精神上的塌。

\begin{shiyu-note}
\label{note:chap10}

\textbf{AI越强，人越不能退后。}

这不是反对AI。这是反对把判断权交给AI。

AI可以成为第二双手。但不能成为替你活着的理由。AI可以搭出桥的形状。但人要决定过不过桥、往哪里过、为什么过。

教育最应该守住的，不只是知识的交付，更是\textbf{判断}——判断什么值得学、什么值得做、什么值得守住。我们无法替未来的AI能力画一条永恒边界，但可以决定：人的判断不能因为工具便利而退出。身体经验也会让判断与后果、他人和具体世界发生联系，这正是工程教育不能轻易外包的部分。
\end{shiyu-note}

\vspace{1cm}

\begin{decision-point}
\label{decision:chap10}

\noindent\textbf{这一章，我们看到了判断权为什么是教育的底线。}

\vspace{0.3cm}

\noindent 下一次，当你用AI生成一份教案、一份报告、一张图的时候——停下来。在点``使用''之前，问自己三个问题：

\noindent 1. 这个东西，哪里可能不对？
\noindent 2. 如果我的学生/同事问我``为什么选这个方案''，我能解释吗？
\noindent 3. 这个东西里，有没有一句话、一个选择、一个细节——是我自己做的？

\noindent 如果第三个问题的答案是``没有''——这份东西不是你的。重做。
\end{decision-point}

\vspace{0.5cm}

\begin{flushright}
{\small\color{muted} 下一章：教师是``改稿机器''吗？——AI时代的教育伦理。}
\end{flushright}
